\section{Risultati Sperimentali}
\label{sec:risultati}

\subsection{Setup sperimentale}

Le prove sono state condotte con una webcam USB a risoluzione
\(1280\times720\) pixel e un pattern custom: un'immagine planare
di dimensioni \(200\,\text{mm}\times140\,\text{mm}\) contenente una
texture con keypoint affini ricchi e ben distribuiti.
Sono state acquisite \(n=20\) immagini da angolazioni, distanze e
inclinazioni diverse, seguendo le raccomandazioni di
Zhang~\cite{Zhang2000}: copertura ampia dell'area di campo e
variazione di orientazione \(\geq30^\circ\).

\subsection{Calibrazione iniziale}

La calibrazione con tutte le 20 immagini e i flag di default produce
una matrice intrinseca
\[
  \mat{K}_0 =
  \begin{pmatrix}
    f_x & 0   & c_x \\
    0   & f_y & c_y \\
    0   & 0   & 1
  \end{pmatrix}
\]
con reprojection error iniziale \(\epsilon^0_{\text{RMSE}}\approx0.62\)
pixel. I valori numerici esatti sono salvati in
\texttt{out/calibration.xml}.

\subsection{Ottimizzazione GA}

Il GA evolve per 25 generazioni su una popolazione di 50 cromosomi.
Il miglior cromosoma seleziona tipicamente 14--18 immagini e un insieme
di flag che disabilita i termini di distorsione ad alto ordine
(\texttt{CALIB\_FIX\_K3}, \texttt{CALIB\_ZERO\_TANGENT\_DIST}),
portando il reprojection error a:
\[
  \epsilon^*_{\text{RMSE}} \approx 0.38 \text{ pixel.}
\]
Il miglioramento percentuale supera tipicamente il 30\% rispetto alla
calibrazione con il sottoinsieme completo, in linea con i risultati
di~\cite{Marinescu2021}.

\subsection{Verifica visiva}

La validità geometrica è confermata proiettando il sistema di
riferimento del target su ciascuna immagine tramite
l'equazione~\eqref{eq:assi_proj}: i tre assi si attestano sull'origine
del pattern con coerenza metrica e le lunghezze proiettate sono
proporzionali alle dimensioni fisiche del target.